\section{EBS and Local Instance Store}

\subsection{EBS}
\begin{itemize}
    \item Network drive you attach to one instance only.
    \item Linked to a specific availability zone (transfer: snapshot -> restore)
    \item Volumes can be resized
    \item Make sure to choose an instance type tha tis EBS optimised to enjoy maximum throughput
\end{itemize}
%TODO add diagram

\subsection{EBS Volume Types}
\begin{itemize}
    \item EBS Volumes come in 6 types
    \begin{itemize}
        \item gp2/gp3 (SSD): General purpose SSD volume that balances price and performance for a wide variety of workloads
        \item io1/io2 Block Express - Highest-performance SSD volume for mission critical low latency or high throughput workloads
        \item st1 (HDD) low cost HDD volume designed for frequently access, throughput intesive workloads
        \item sc1 (HDD): Lowest cost HDD volume designed for less frequently accessed workloads
    \end{itemize}
    \item EBS volumes are characterised in Size / throughput / IOPS (I/O Ops Per Sec)
    --- See AWS Docs when required
    \item Only gp2/gp3 and io1/io2 can be used as boot volumes
\end{itemize}

\subsection{EBS Snapshots}
\begin{itemize}
    \item Incremental - only backup changed blocks
    \item EBS backups use IO and you shouldn't run them while your application is handling a lot of traffic
    \item Snapshots will be stored in S3 (you won't directly see them)
    \item Not necessary to detach volume to do snapshot, but recommended
    \item Can copy snapshots across region (for DR)
    \item Can make Image (AMI) from snapshot
    \item EBS volumes restored by snapshots need to be pre-warmed (use the Fast Snapshot Restore FSR feature or fio/dd command to read the entire volume)
\end{itemize}

\subsection{Amazon Data Lifecycle Manager}
\begin{itemize}
    \item Automate the creation, retention and deletion of EBS snapshots and EBS-backed AMIs
    \item Schedule backups, cross-account snapshot copies delete outdated backups
    \item Uses resource tags to identify the resources (EC2 instances, EBS volumes)
    \item Can't be used to manage snapshots/AMIs created outside Data Life Cycle manager
    \item Can't be used to manage instance-store backed AMIs
\end{itemize}

\subsection{Amazon Data Lifecycle Manger vs AWS Backup}
\begin{itemize}
    \item Use Data Lifecycle Manager - when you want to automate the creations, retention and deletion of EBS snapshots
    \item Use AWS backup - to manage and monitor backups across the AWS services you use, including EBS volumes from single place
\end{itemize}

\subsection{EBS Encryption - Account level setting}
\begin{itemize}
    \item New Amazon EBS volumes aren't encrypted by default
    \item There's an account-level setting to encrypt automatically EBS volumes and Snapshots
    \item The setting needs to be enabled on a per-region basis
\end{itemize}

\subsection{EBS Multi-Attach - io1/io2 family}
\begin{itemize}
    \item Attach the same EBS volume to multiple EC2 instances in the same AZ
    \item Each instance has full read and write permissions to the volume
    \item Use case:
    \begin{itemize}
        \item Achieve higher application availability in clustered in clustered Linux applications (ex: Teradata)
        \item Applications must manage concurrent write operations
        \item Must use a file system that's cluster-aware (not XFS, EX4, etc ...)
    \end{itemize}
\end{itemize}

\subsection{Local EC2 instance store}
\begin{itemize}
    \item Physical disk attached to the physical server where your EC2 is
    \item Very High IOPS (because its physical)
    \item Disks up to 7.5 TiB (can change over time), stripped to reach 60 TiB (can change over time)
    \item Block Storage (just like EBS)
    \item Cannot be increased in size
    \item Risk of data loss if hardware fails
\end{itemize}

\subsection{Instance Store vs EBS}
\begin{itemize}
    \item Instance store is physically attached to the machine (ephemeral storage)
    \item EBS is a network drive (persistent)
    \item Pros
    \begin{itemize}
        \item Better I/O performance (EBS gp2 has a max IOPS of 16000, io1 of 6400, io2 Block Express of 256000)
        \item Good for buffer / cache / scratch data / temporary content
        \item Data survives reboots
    \end{itemize}
    \item Cons
    \begin{itemize}
        \item On stop or termination , the instance store is lost
        \item You can't resize the instance store
        \item Backups must be operated by the user
    \end{itemize}
\end{itemize}


\section{Amazon EFS}
\begin{itemize}
    \item Managed NFS (network file system) that can be mounted on many EC2
    \item EFS works with EC2 instances in multi-AZ and on-premises (DX and VPN)
    \item Highly available, scalable, expensive (3x gp2), pay per GB used
\end{itemize}
%TODO add diagram

\subsection{EFS - Elastic File System}
\begin{itemize}
    \item Use cases: content management, web serving, data sharing, WordPress
    \item Compatible with Linux based AMI (not Windows), POSIX-compliant
    \item Uses NFSv4.1 protocol
    \item Uses security group to control access to EFS
    \item Encryption at rest using KMS
    \item POSIX file system (roughly equal to Linux) that has a standard file API
    \item File system scala automatically, pay per-use, no capcity planning
\end{itemize}

\subsection{EFS - Performance and Storage Classes}

\subsubsection{EFS Scale}
\begin{itemize}
    \item 1000s of concurrent EFS clients, 10 GB+ /s throughput
    \item Grow to Petabyte-scala network file system, automatically
    \item Performance Mode (set as EFS creation time)
    \item General Purpose (default) - latency-sensitive use cases (web server, CMS, etc.....)
    \item Max I/O - higher latency, throughput, highly parallel (big data, media processing)
    \item Throughput Mode
    \begin{itemize}
        \item Bursting - 1 TB = 50 MiB/s + burst of up to 100MiB/s
        \item Provisioned - set your throughput regardless of storage size ex: 1GiB/s for 1 TB storage
        \item Elastic - automatically scales throughput up or down based on your workloads
        \item Up to 3GiB/s for reads and 1GiB/s for writes
        - Used for unpredictable workloads
    \end{itemize}
\end{itemize}

\subsection{EFS - Storage Classses}
\begin{itemize}
    \item Storage Tiers (Lifecycle management feature - move file after N days)
    \begin{itemize}
        \item Standard: for frequently accessed files
        \item Infrequent access (EFS-1A): cost to retrieve files, lower price to store
        \item Archive: rarely accessed data (few times each year), 50\% cheaper
        \item Implement lifecycle policies to move files between storage tiers
    \end{itemize}
    \item Availability and durability
    \begin{itemize}
        \item Standard: Multi-AZ great for prod
        \item One Zone: One AZ great for dev, backup enabled by default, compatible with IA (EFS One Zone-IA)
    \end{itemize}
    \item Over 90\% in cost savings
\end{itemize}
%TODO add diagram

\subsection{EFS - On-premises and VPC peering}
%TODO add diagram

\subsection{EFS - Access Points}
\begin{itemize}
    \item Easily manage applications access to NFS environments
    \item Enforce a POSIX user and group to use when accessing the file system
    \item Restrict access to a directory within the file system and optionally specify a different root directory
    \item Can restrict access from NFS clients using IAM policies
\end{itemize}
%TODO add diagram

\subsection{EFS - File System Policies}
\begin{itemize}
    \item Resource - based policy to control access to EFS file systems (same as S3 bucket policy)
    \item By default, it grants full access to all clients
\end{itemize}
%TODO add JSON example

\subsection{EFS - Cross Region Replication}
\begin{itemize}
    \item Replicate objects in an EFS file system to another AWS regions
    \item Setup for new or existing EFS file systems
    \item Provides RPO and RTO of minutes
    \item Doesn't affect the provisioned throughput of the EFS file system
    \item Use cases: meet you compliance and business continuity goals
\end{itemize}
%TODO add JSON example


\section{Amazon S3}

\subsection{S3 - Overview}
\begin{itemize}
    \item Object storage, serverless, unlimited storage, pay-as-you-go
    \item Good to store static content (image, video files)
    \item Access object by key, no indexing facility
    \item Not a filesystem, cannot be mounted natively on EC2

    \item Anti patterns:
    \begin{itemize}
        \item Lots of small files
        \item POSIX file system (use EFS instead), file locks
        \item Search features, queries, rapidly changing data
        \item Website with dynamic content
    \end{itemize}
\end{itemize}

\subsection{S3 Storage Classes Comparison}
%TODO add table / see table

- You can transition object between tiers (or delete) using S3 Lifecycle policies
--> See URL

\subsection{S3 - Replication (Versioning Enabled)}
\begin{itemize}
    \item Cross Region Replication (CRR)
    \item Same Region Replication (SRR)
    \item Combine with Lifecycle rules
    \item Helpful to reduce latency, disaster recovery, security
    \item S3 Replication Time Control (S3 RTC)
    \begin{itemize}
        \item Replicates most objects that you upload to Amazon S3 in seconds, and 99.99\% of those objets within 15 minutes
        \item Helpful for compliance, DR etc
    \end{itemize}
\end{itemize}
%TODO add diagram

\subsection{S3 Event Notifications}
%TODO add diagram
\begin{itemize}
    \item S3:ObjectCreated, S3:ObjectRemoved, S3:ObjectRestore, S3:Replication,......
    \item Object name filtering possible (*.jpg)
    \item Use case: generate thumbnails of images uploaded to S3
    \item Can create as many "S3 events" as desired
    \item S3 event notifications typically deliver events in seconds but can sometimes take a minute or longer
\end{itemize}

\subsection{S3 Event Notification with Amazon EventBridge}
%TODO add diagram
\begin{itemize}
    \item Advanced filtering options with JSON rules (metadata, object size, name.......)
    \item Multiple Destinations - ex Step Functions, Kinesis Streams / Firehose ......
    \item EventBridge Capabilities - Archive, Replay Events, Reliable Delivery
\end{itemize}

\subsection{S3 - Baseline Performance}
\begin{itemize}
    \item Amazon S3 automatically scales to high request rates, latency 100-200ms
    \item Your application can achieve at least 3500 PUT/COPY/POST/DELETE or 5,500 GET/HEAD reqeusts per second per prefix in a bucket
    \item There are limits to the number of prefixes in a bucket
    \item Example (object path => prefix)
    \begin{itemize}
        \item bucket/folder1/sub1/file => /folder1/sub1/
        \item bucket/folder1/sub2/file => /folder1/sub2/
        \item bucket/1/file => /1/
        \item bucket/2/file => /2/
    \end{itemize}
    \item If you spread reads across all four prefixes evenly, you can achieve 22,000 requests per second for GET and HEAD
\end{itemize}

\subsection{S3 Performance}
\begin{itemize}
    \item Multi-Part upload:
    \begin{itemize}
        \item recommended for files > 100MB, must use for files > 5GB
        \item can help parallelize uploads (speed up transfer)
    \end{itemize}
%TODO diagram
    \item S3 Transfer Acceleration
    \begin{itemize}
        \item Increase transfer speed by transferring file to an AWS edge location which will forward the data to the S3 bucket in the target region
        \item Compatible with multi-part upload
    \end{itemize}
%TODO diagram
\end{itemize}

\subsection{S3 Performance - S3 Byte-Range Fetches}
\begin{itemize}
    \item Parallelize GETs by requesting specific byte ranges
    \item Better resilience in case of failures
\end{itemize}
\underline{Can be used to speed up downloads}
%TODO add diagram

\underline{Can be used to retrieve only partial data (for example the head of a file)}
%TODO add diagram

\subsection{S3 Multi-Part - Remove Incomplete Parts}
%TODO add diagram


\section{Amazon S3 - Storage Class Analysis}
\begin{itemize}
    \item May be seen as "storage class analysis" at the exam
    \item Help you decide when to transition objects to the right storage class
    \item Recommendations for Standard and Standard IA
    \begin{itemize}
        \item Does NOT work for One-Zone IA or Glacier
    \end{itemize}
    \begin{itemize}
        \item Report is updated daily
        \item 24 to 48 hours to start seeing data analysis
        \item Visualise data in Amazon QuickSight
        \item Good first step to put together Lifecycle Rules (or improve them!)
    \end{itemize}
\end{itemize}


\section{Amazon S3 - Storage Lens}
\begin{itemize}
    \item Understand, analyse and optimise storage across entire AWS Organisation
    \item Discover anomalies, identify cost efficiencies and apply data protection best practices across entire AWS Organisation (30 days usage and activity metrics)
    \item Aggregate data for Organisations, specific accounts, regions, buckets or prefixes
    \item Default dashboard or create your own dashboards
    \item Can be configured to export metrics daily to an S3 bucket (CSV, Parquet)
\end{itemize}
%TODO diagram

\subsection{Storage Lens - Default Dashboard}
\begin{itemize}
    \item Visualise summarised insights and trends for both free and advanced metrics
    \item Default dashboard shows Multi-Region and Multi-Account data
    \item Preconfigured by Amazon S3
    \item Can't be deleted, but can be disabled
\end{itemize}
%TODO see screenshots

\subsection{Storage Lens - Metrics}
\begin{itemize}
    \item Summary Metrics
    \begin{itemize}
        \item General insights about your S3 storage
        \item StorageBytes, ObjectCount.....
        \item Use cases: identify the fastest-growing (or no used) buckets and prefixes
    \end{itemize}
    \item Cost-Optimised Metrics
    \begin{itemize}
        \item Provide insights to managed and optimise your storage costs
        \item NonCurrentVersionStorageBytes, IncompleteMultipartUploadStorageBytes...
        \item Use cases: identify buckets with incomplete multipart uploaded older than 7 days, Identify which objects could be transitioned to lower-cost storage class
    \end{itemize}
    \item Data-Protection Metrics
    \begin{itemize}
        \item Provide insights for data protection features
        \item VersioningEnabledBucketCount, MFADeleteEnabledBucketCount, SSEKMSEnabledBucketCount, CrossRegionReplicationRuleCount...
        \item Use cases: identify buckets that aren't following data-protection best practices
    \end{itemize}
    \item Access-management Metrics
    \begin{itemize}
        \item Provide insights for S3 Object Ownership
        \item ObjectOwnershipBucketOwnerEnforcedBucketCount...
        \item Use cases: identify which Object Ownership settings your buckets use
    \end{itemize}
    \item Event Metrics
    \begin{itemize}
        \item Provide insights for S3 Event Notifications
        \item EventNotificationEnabledBucketCount (identify which buckets have S3 Event Notifications)
    \end{itemize}
    \item Performance Metrics
    \begin{itemize}
        \item Provide insights for S3 Transfer Acceleration
        \item TransferAccelerationEnabledBucketCount (identify which buckets have S3 Transfer Acceleration enabled)
    \end{itemize}
    \item Activity Metrics
    \begin{itemize}
        \item Provide insights about how your storage is requested
        \item AllRequests, GetRequests, PutRequests, ListRequests, BytesDownloaded
    \end{itemize}
    \item Detailed Status Code Metrics
    \begin{itemize}
        \item Provide insights for HTTP status codes
        \item 200OKStatusCount, 403ForbiddenErrorCount, 404NotFoundErrorCount...
    \end{itemize}
\end{itemize}

\subsection{Storage Lens - Free vs Paid}
\begin{itemize}
    \item Free Metrics
    \begin{itemize}
        \item Automatically available for all customers
        \item Contain around 28 usage metrics
        \item Data is available for queries for 14 days
    \end{itemize}
    \item Advanced Metrics and Recommendations
    \begin{itemize}
        \item Additional paid metrics and features
        \item Advanced Metrics - Activity, Advanced Cost Optimisation, Advanced Data Protection, Status code
        \item CloudWatch Publishing - Access Metrics in CloudWatch without additional charges
        \item Prefix-Aggregation - Collect metrics at the prefix level
        \item Data is available for queries for 15 months
    \end{itemize}
\end{itemize}


\section{S3 Solution Architecture}

%TODO diagram

\subsection{S3 Solution Architecture Indexing objects in Dyanmo DB}
%TODO add diagram
---- Notes for DynamoDB table

\subsubsection{API for object metadata}
\begin{itemize}
    \item Search by date
    \item Total storage used by a customer
    \item List of all objects with certain attributes
    \item Find all objects uploaded within a date range
\end{itemize}

\subsection{Solution Architecture on AWS Dynamic vs Static Content}
%TODO add diagram


\section{Amazon FSx}

\subsection{Amazon FSx - Overview}
\begin{itemize}
    \item Launch 3rd party high-performance file systems on AWS
    \item Fully managed service
\end{itemize}
%TODO logos

\subsection{Amazon FSx for Windows (File Server)}
\begin{itemize}
    \item FSx for Windows is fully managed Windows file system share drive
    \item Supports SMB protocol and Windows NTFS
    \item Microsoft Active Directory integration, ACLs, user quotas
    \item Can be mounted on Linux EC2 instances
    \item Supports Microsoft's Distributed File System (DFS) Namespaces (group files across multiple FS)
    \item Scale up to 10s of GB/s, millions of IOPS, 100s PB of data
    \item Storage Options:
    \begin{itemize}
        \item- SSD - latency sensitive workloads (databases, media processing, data analytics....)
        \item- HDD - broad spectrum of workloads (home directory, CMS,....)
    \end{itemize}
    \item Can be accessed from your on-premises infrastructure (VPN or Direct Connect)
    \item Can be configured to be Multi-AZ (high availability)
    \item Data is backed-up daily to S3
\end{itemize}

\subsection{Amazon FSx for Lustre}
\begin{itemize}
    \item Lustre is a type of parallel distributed file system, for large-scale computing
    \item The name Lustre is derived from "Linux" and "cluster"
    \item Machine Learning, High Performance Computing (HPC)
    \item Video Processing, Financial Modeling, Electronic Design Automation
    \item Scales up to 100s GB/s, millions IOPS, sub-ms latencies
    \item Storage Options:
    \begin{itemize}
        \item SSD - low-latency, IOPS intensive workloads, small and random file operations
        \item HDD - throughput-intensive workloads, large and sequential file operations
    \end{itemize}
    \item Seamless integrations with S3
    \begin{itemize}
        \item Can "read S3" as a file system (through FSx)
        \item Can write the output of the computations back to S3 (through FSx)
    \end{itemize}
    \item Can be used from on-premises servers (VPN or Direct Connect)
\end{itemize}

\subsection{FSx Lustre - File System Deployment Options}
%TODO add diagram
\begin{itemize}
    \item Scratch File System
    \begin{itemize}
        \item Temporary storage
        \item Data is not replicated (doesn't persist if file server fails)
        \item High burst (6x faster, 200MBps per TiB)
        \item Usage: short-term processing, optimised costs
    \end{itemize}
    \item Persistent File System
    \begin{itemize}
        \item Long-term storage
        \item Data is replicated within same AZ
        \item Replace failed files within minutes
        \item Usage: long-term processing, sensitive data
    \end{itemize}
\end{itemize}

\subsection{Amazon FSx for NetApp ONTAP}
\begin{itemize}
    \item Managed NetApp ONTAP on AWS
    \item File System compatible with NFS, SMB, iSCSI protocol
    \item Move workloads running on ONTAP or NAS to AWS
    \item Work with:
    \begin{itemize}
        \item Linux
        \item Windows
        \item MacOS
        \item VMWare Cloud on AWS
        \item Amazon Workspaces and AppStream 2.0
        \item Amazon EC2, ECS and EKS
    \end{itemize}
    \item Storage shrinks or grows automatically
    \item Snapshots, replication, low-cost, compression and data de-duplication
    \item Point-in-time instantaneous cloning (helpful for testing new workloads)
\end{itemize}
%TODO diagram

\subsection{Amazon FSx for OpenZFS}
\begin{itemize}
    \item Managed OpenZFS file system on AWS
    \item File System compatible with NFS (v3, v4, v4.1, v4.2)
    \item Move workloads running on ZFS to AWS
    \item Works with:
    \begin{itemize}
        \item Linux
        \item Windows
        \item MacOS
        \item VMware Cloud on AWS
        \item Amazon Workspaces and AppStream 2.0
        \item Amazon EC2, ECS and EKS
    \end{itemize}
    \item Up to 1,000,000 IOPS with < 0.5ms latency
    \item Snapshots, compression and low-cost
    \item Point-in-time instantaneous cloning (helpful for testing new workloads)
\end{itemize}
%TODO add diagram


\section{Amazon FSx - Solution Architecture}
%TODO add diagram of FSx - Solution Architecture

\subsection{FSx - Solution Architecture Decrese FSx Volume Size}
\begin{itemize}
    \item If you take a backup, you can only restore to a same size
    \item You can only increase the amount of storage capacity for a file system; you cannot decrease storage capacity
    \item Instead, create a new FSx (smaller), use DataSync to sync data and then migrate your app over
\end{itemize}
%TODO add diagram

\subsection{FSx for Lustre - Data Lazy Loading}
%TODO add diagram
\begin{itemize}
    \item Any data processing job on Lustre with S3 as an input data source can started withouth Lustre doing a full download of the data set first
    \item Data is lazy loaded: only the data this is actually processed is loaded, meaning you can decrease your costs and latency
    \item Data is also loaded only once, therefore you reduce your requests on Amazon S3
\end{itemize}


\section{AWS DataSync}
\begin{itemize}
    \item Move large amount of data to and from
    \begin{itemize}
        \item On-premises / other cloud to AWS (NFS, SMB, HDFS, S3 API...) - needs agent
        \item AWS to AWS (different storage services) - no agent needed
    \end{itemize}
    \item Can synchronise to:
    \begin{itemize}
        \item Amazon S3 (any storage classes, including glacier)
        \item Amazon EFS
        \item Amazon FSx (Windows, Lustre, NetApp, OpenZFS)
    \end{itemize}
    \item Replication tasks can be scheduled hourly, daily, weekly
    \item File permissions and metadata are preserved (NFS POSIX, SMB...)
    \item One agent task can use 10Gbps, can setup a bandwidth limit
\end{itemize}


\section{AWS DataSync - Solution Architecture}

\subsection{AWS DataSync NFS / SMB to AWS (S3, EFS, FSx)}
%TODO add diagram

\subsection{AWS Datasync Transfer between AWS storage services}
%TODO add diagram

\subsection{AWS DataSync Private VIF through Direct Connect}
%TODO add diagram


\section{AWS Data Exchange}
\begin{itemize}
    \item Find, subscribe to and use third - party data in the cloud
    \begin{itemize}
        \item Reuters, who curate data from over 2.2 million unique news stories per year in multiple languages
        \item Change healthcare, who process and anonymize more that 14 billion healthcare transactions and \$1 trillion in claims annually
        \item Dun and Bradsreet, who maintain a database of more than 330 million global business records;
        \item Foursquare, whose location data is derived from 220 million unique consumers and includes more 60 million global commercial venues
    \end{itemize}
    \item Once subscribed to a data product, you can use the AWS data Exchange API to load data directly into Amazon S3 and then analyze it with wide variety to AWS analytics and machine learning serivces
\end{itemize}

\subsection{AWS Data Exchange - Other Products}
\begin{itemize}
    \item AWS Data Exchange for Redshift
    \begin{itemize}
        \item Find and subscribe to third-party data in AWS Data Exchange that you can query in an Amazon Redshift data warehouse in minutes
        \item Easily license your data in Amazon Redshift through AWS Data Exchange
    \end{itemize}
    \item AWS Data Exchange for APIs
    \begin{itemize}
        \item Find and subscribe to third-party APIs with a consistent access using AWS SDKs
        \item Consistent AWS-native authentication and governance
    \end{itemize}
\end{itemize}


\section{AWS Transfer Family}
\item A fully-managed service for file transfers into and out of Amazon S3 or Amazon EFS using the FTP protocol
\item Supported Protocols
\begin{itemize}
    \item AWS Transfer for FTP (File Transfer Protocol (FTP))
    \item AWS Transfer for FTPS (File Transfer Protocol over SSL (FTPS))
    \item AWS Transfer for SFTP (Secure File Transfer Protocol (SFTP))
\end{itemize}
\item Managed infrastructure, Scalable, Reliable, Highly Available (multi-AZ)
\item Pay per provisioned endpoint per hour + data transfers in GB
\item Store and manage users' credentials within the service
\item Integrate with existing authentication systems (Microsoft Active Directory, LDAP, Okta, Amazon Cognito, custom)
\item Usage: sharing files, public datasets, CRM, ERP,.....

%TODO add diagram

\subsection{AWS Transfer Family - Endpoint Types}

\subsubsection{Public Endpoint}
%TODO add diagram
\begin{itemize}
    \item IPs managed by AWS (subject to change, use DNS names)
    \item Can't setup allow lists by source IP addresses
\end{itemize}

\subsubsection{VPC Endpoint with Internal Access}
%TODO add diagram
\begin{itemize}
    \item Static private IPs
    \item Setup allow lists (SGs and NACLs)
\end{itemize}

\subsubsection{VPC Endpoint with Internet-facing Access}
%TODO add diagram
\begin{itemize}
    \item Static private IPs
    \item Static public IPs (EIPs)
    \item Setup security groups
\end{itemize}


\section{AWS Storage Services Price Comparison}
%TODO add diagram